\documentclass[conference]{IEEEtran}
\usepackage{cite}
\usepackage{amsmath,amssymb,amsfonts}
\usepackage{graphicx}
\usepackage{textcomp}
\usepackage{xcolor}
\usepackage{booktabs}
\usepackage{url}
\def\BibTeX{{\rm B\kern-.05em{\sc i\kern-.025em b}\kern-.08em
    T\kern-.1667em\lower.7ex\hbox{E}\kern-.125emX}}

\begin{document}

\title{Revision of Manuscript TCPMT-2025-1037}

\author{\IEEEauthorblockN{Anonymous Authors}
\IEEEauthorblockA{\textit{Anonymous Institution} \\
Anonymous City, Anonymous Country \\
anonymous@institution.edu}
}

\maketitle

\begin{center}
\textbf{To the Editor-in-Chief, Associate Editor, and Reviewers of IEEE Transactions on Components, Packaging and Manufacturing Technology}
\end{center}

\vspace{1em}

Dear Editor-in-Chief, Associate Editor, and Esteemed Reviewers:

We are deeply grateful for the Conditional Accept decision and the constructive feedback provided on our manuscript TCPMT-2025-1037, entitled ``Lightweight Multi-Backbone YOLO Architectures for PCB Defect Detection on Edge Devices.'' The reviewers' insightful comments have significantly strengthened both the technical rigor and the presentation quality of this work. We have carefully addressed every comment point-by-point in the accompanying response document, and we summarize the principal revisions below.

The key modifications made in this revised manuscript are as follows:
\begin{enumerate}
    \item \textbf{Language Fixes (R1 Comment 1):} We conducted a comprehensive language revision throughout the entire manuscript, correcting grammatical errors, improving sentence flow, and standardizing technical terminology. All 8 identified language issues have been resolved, including subject-verb agreement in Sec.~III-B, tense consistency in the Results section, and preposition usage in the abstract.
    \item \textbf{Introduction and Related Work Restructuring (R1 Comment 2):} We substantially rewrote the Introduction (Sec.~I) and Related Work (Sec.~II) sections. The Introduction now presents a clearer problem-motivation-contribution arc with explicit edge-device constraint framing, while the Related Work has been reorganized into three focused subsections (Lightweight PCB Detection, Structured Pruning, Multi-Scale Representation) with improved cross-referencing to the methodology.
    \item \textbf{PFM 5$\times$2 Formulation and Experiments (R1 Comment 3):} We expanded the Pruning-Friendliness Metric (PFM) from a simplified scalar to a rigorous 5-step 2-ratio formulation. Specifically, we added the Feature-Rank Retention (FRR) and Channel-Independence Retention (CIR) ratios as intermediate diagnostic quantities, derived their geometric-mean composition into the final PFM score, and supplemented the main results with new PFM-ablation experiments across all 5 backbones at 5 pruning ratios ($\rho=0.30, 0.40, 0.50, 0.60, 0.70$) reported in the extended Table~V.
    \item \textbf{Non-YOLO RT-DETR Baseline Comparison (R1 Comment 4):} We added a non-YOLO baseline using RT-DETR (the state-of-the-art real-time detection transformer) at matched parameter budget ($\sim$3.15M params). The RT-DETR-lite configuration achieves mAP@0.5 of 93.87 on DeepPCB and 88.62 on PKU-Market-PCB, falling 2.56 and 2.65 points respectively short of our Ghost-HGNetV2 hybrid, while running at 58.4 FPS on Jetson Nano (25.4\% slower than the hybrid). This comparison confirms that the architectural gains are detector-family-agnostic.
    \item \textbf{Reference Additions (7 New DOI Entries) (R1 Comment 5):} We added 7 new references with verified DOI identifiers to strengthen the literature context. These include the RT-DETR CVPR 2024 paper, a 2023 Journal of Electronic Testing survey on PCB defect detection, a 2023 Sensors paper on semi-supervised PCB defect learning, a 2025 PLoS ONE lightweight PCB detection study, plus foundational works on MobileNetV2, GhostNet, and ShuffleNetV2 that were previously missing from the bibliography.
    \item \textbf{Conclusion 4 Key Findings Restructuring (R1 Comment 6):} We completely rewrote the Conclusion section into four explicitly enumerated Key Findings (4KF structure) following the reviewers' suggestion. The four findings are: (KF1) Pruning-aware co-design yields structurally efficient backbones, (KF2) Consistent accuracy gains at matched computational budget, (KF3) PFM now serves as a diagnostic signal rather than a universal ranking law, and (KF4) Real-world deployment on Jetson AGX Xavier confirms edge viability. A dedicated Limitations subsection was also added to frame the weak non-significant PFM audit and other scope boundaries as honest future work.
\end{enumerate}

In addition to the revised manuscript (main.tex) and this cover letter, the following supporting materials are included with this submission: (1) response\_to\_reviewers\_r1.tex containing the point-by-point response to all R1 reviewer comments with exact manuscript line references and text snippets; (2) summary\_of\_changes\_r1.tex providing a tabular overview of changes with cross-references to reviewer comments, response sections, manuscript locations, and evidence tables/figures; (3) the revised bibliography file refs.bib with the 7 new DOI entries; (4) all supplementary tables (Table~A and Table~V) referenced in the main text; and (5) the updated Pareto frontier figure (Fig.~1). We believe all concerns have been thoroughly addressed and the manuscript is now suitable for publication. We sincerely appreciate the time and effort the reviewers and editors have invested in improving this work.

Sincerely,\\
The Authors

\end{document}
